\section{Berechnung der Beiträge für den Neuzugang}\label{sec:beitraege}

Die Beitragsberechnung folgt den Formeln des Abschnitts~A der Anlage~1
zur KVAV (\S\,10 Abs.\,1 KVAV) in der nachfolgend dokumentierten
rekursiven Darstellung. Da der Tarif für das Neugeschäft geschlossen
ist (Abschnitt~\ref{sec:tarifdarstellung}), dienen die
Neuzugangsbeiträge als Tarifbeiträge der Bestandsführung, der
Beitragsanpassung (Abschnitt~\ref{sec:bap}) und individuellen
Tarifänderungen (Abschnitt~\ref{sec:tarifaenderung}). Das
Höchsteintrittsalter des Tarifs betrug 75~Jahre; Tarifbeiträge für
höhere Alter werden zu tariftechnischen Zwecken geführt.

\subsection{Versicherungstechnische Grundwerte}\label{sec:grundwerte}

Die Barwerte werden je Beobachtungseinheit (Männer, Frauen) rekursiv
aus den Rechnungsgrundlagen des Abschnitts~\ref{sec:rechnungsgrundlagen}
berechnet. Der \textbf{Leistungsbarwert} $A_x$ (in Einheiten des
Grundkopfschadens, d.\,h. in Profileinheiten) und der
\textbf{Rentenbarwert} $a_x$ (Barwert einer lebenslang vorschüssig
gezahlten Rente der Höhe 1 p.\,a.) erfüllen die Rekursionen
\begin{align}
A_{110} &= k_{110}, &
A_x &= k_x + DQ_x\cdot A_{x+1}, \qquad x = 109,\dots,21,
\label{eq:rekA}\\
a_{110} &= 1, &
a_x &= 1 + DQ_x\cdot a_{x+1}, \qquad x = 109,\dots,21,
\label{eq:reka}
\end{align}
wobei $k_x$ den Profilwert gemäß Gleichung~\eqref{eq:profil} und
$DQ_x$ die diskontierte Übergangswahrscheinlichkeit gemäß
Gleichung~\eqref{eq:dq} bezeichnen. Leistungen und Beiträge werden
rechnungsmäßig jeweils zu Jahresbeginn fällig (vorschüssige
Betrachtung). In Summenschreibweise gilt gleichwertig
\begin{equation}
A_x = \sum_{t=0}^{110-x} v^t \cdot {}_{t}p_x \cdot k_{x+t},
\qquad
a_x = \sum_{t=0}^{110-x} v^t \cdot {}_{t}p_x,
\qquad
{}_{t}p_x = \prod_{s=0}^{t-1}\bigl(1-(q+w)_{x+s}\bigr),
\label{eq:barwertsumme}
\end{equation}
wobei ${}_{t}p_x$ die $t$-jährige Verbleibewahrscheinlichkeit einer
$x$-jährigen Person bezeichnet. Der Leistungsbarwert in EUR ergibt sich
durch Multiplikation mit dem Grundkopfschaden: $G\cdot A_x$. Die
Barwerte enthält Anlage~\ref{sec:anlage311}; die
Rundungskonventionen des Abschnitts~\ref{sec:rundung} gelten
schrittweise in den Rekursionen.

\subsection{Jahresnettobeiträge}\label{sec:nettobeitraege}

Der Jahresnettobeitrag ergibt sich für Erwachsene nach dem
Äquivalenzprinizip (Barwert der Leistungen = Barwert der Beiträge) zu
\begin{equation}
P_x = \frac{G\cdot A_x}{a_x},
\qquad x = 21,\dots,110,
\label{eq:netto}
\end{equation}
wobei $G$ den Grundkopfschaden, $A_x$ den Leistungsbarwert in
Profileinheiten und $a_x$ den Rentenbarwert bezeichnen ($x$ ist das
Eintrittsalter; der Nettobeitrag bleibt bei unveränderten
Rechnungsgrundlagen lebenslang konstant).

Kinder und Jugendliche zahlen die natürliche Prämie (\S\,10 Abs.\,4
KVAV); ihr Jahresnettobeitrag entspricht dem mittleren Kopfschaden der
Altersgruppe gemäß Gleichung~\eqref{eq:kkinder}. Die Jugendlichen werden
dabei geschlechtsabhängig unterschieden:
\begin{equation}
P^{\mathrm{Ki}} = \bar{K}_{\mathrm{Ki}},
\qquad
P^{\mathrm{Ju},M} = \bar{K}^{M}_{\mathrm{Ju}},
\qquad
P^{\mathrm{Ju},F} = \bar{K}^{F}_{\mathrm{Ju}}.
\label{eq:pkinder}
\end{equation}

\subsection{Tarifliche monatliche Beitragsraten}\label{sec:beitragsraten}

Der tarifliche monatliche Bruttobeitrag der Erwachsenen
(Einzelversicherung) beträgt
\begin{equation}
b_x = \frac{G\cdot A_x + \gamma_x\cdot a_x}
           {12\cdot\bigl(1-(\sigma+\Omega_x)\bigr)\cdot a_x - \alpha_x},
\qquad x = 21,\dots,110,
\label{eq:brutto}
\end{equation}
wobei
\begin{itemize}\itemsep1pt
  \item $G\cdot A_x$ der Leistungsbarwert in EUR,
  \item $\gamma_x$ der altersunabhängige Zuschlag gemäß
        Gleichung~\eqref{eq:gammasumme} (Stufe $x \le 64$ bzw.
        $x \ge 65$ nach dem Eintrittsalter),
  \item $\sigma$ der Sicherheitszuschlag gemäß
        Gleichung~\eqref{eq:sigma},
  \item $\Omega_x$ der Zuschlag gemäß Gleichung~\eqref{eq:omegaz},
  \item $\alpha_x$ der Zillmersatz gemäß Gleichung~\eqref{eq:zillmer}
\end{itemize}
ist. Die Formel folgt aus der Äquivalenzgleichung
\begin{equation}
12\,b_x\,a_x = G\,A_x + \gamma_x\,a_x
+ 12\,(\sigma+\Omega_x)\,b_x\,a_x + \alpha_x\,b_x,
\label{eq:aequivalenz}
\end{equation}
also Beitragsbarwert = Leistungsbarwert + Barwert der absoluten Zuschläge
+ Barwert der prozentualen Zuschläge + Zillmerkosten.

Kinder und Jugendliche (keine Alterungsrückstellung, keine Zillmerung;
$\Omega$ in der Stufe $x \le 64$; Jugendliche geschlechtsabhängig):
\begin{equation}
\begin{aligned}
b^{\mathrm{Ki}} &= \frac{\bar{K}_{\mathrm{Ki}} + \gamma^{K}}
{12\,\bigl(1-(\sigma+\Omega_{\le 64})\bigr)}, &
b^{\mathrm{Ju},M} &= \frac{\bar{K}^{M}_{\mathrm{Ju}} + \gamma^{K}}
{12\,\bigl(1-(\sigma+\Omega_{\le 64})\bigr)}, \\[2pt]
b^{\mathrm{Ju},F} &= \frac{\bar{K}^{F}_{\mathrm{Ju}} + \gamma^{K}}
{12\,\bigl(1-(\sigma+\Omega_{\le 64})\bigr)}. &&
\end{aligned}
\label{eq:bkinder}
\end{equation}
Die konkreten Beitragsraten der Kinder sowie der männlichen und
weiblichen Jugendlichen sind in Anhang~\ref{sec:anhang3} ausgewiesen.

Die Beitragsraten werden kaufmännisch auf 0{,}01~EUR gerundet und sind
in Anhang~\ref{sec:anhang3} aufgeführt.
\textbf{Eigenschaften des Beitragsverlaufs:} Die Beiträge sind im
Eintrittsalter monoton nicht fallend, mit zwei systematisch bedingten
Ausnahmen: In den höchsten Altern (ab etwa Alter~90) sind die Beiträge
konstant bis auf Rundungsdifferenzen von wenigen Cent (höchstens
0{,}06~EUR) aus den Rundungskonventionen des
Abschnitts~\ref{sec:rundung}; bei den Männern und Frauen kann der
Beitrag vom Eintrittsalter 64 zum Eintrittsalter 65 geringfügig sinken,
da ab dem Alter 65 die kostenbedingten Zuschläge $\alpha_u$ und die
Standardtarif-Komponente von $\Omega$ entfallen. Beide Effekte stehen
\S\,146 Abs.\,2 VAG nicht entgegen, da sie ausschließlich den
(geschlossenen) Neuzugang betreffen.

\subsection{Gesetzlicher Beitragszuschlag}\label{sec:gz}

Auf den monatlichen Bruttobeitrag wird ab Beginn des Kalenderjahres
nach der Vollendung des 21.~Lebensjahres bis zum Ende des
Kalenderjahres, in dem das 60.~Lebensjahr vollendet wird, der
gesetzliche Beitragszuschlag von 10\,\% gemäß \S\,149 VAG erhoben
(Verwendung: Abschnitt~\ref{sec:zusatzrueckstellung}); dies entspricht
den versicherungstechnischen Altern 22 bis 60. Der monatliche
Zahlbetrag beträgt damit
\begin{equation}
Z_{x_e,m} = \begin{cases}
(1 + 0{,}1 + RZ)\cdot b_{x_e} & 22 \le x_e+m \le 60,\\[2pt]
(1 + RZ)\cdot b_{x_e} & \text{sonst},
\end{cases}
\label{eq:zahlbeitrag}
\end{equation}
wobei $x_e$ das Eintrittsalter, $m$ die abgelaufene Versicherungsdauer
und $RZ$ einen etwaigen individuellen Risikozuschlagsfaktor
(Abschnitt~\ref{sec:risikozuschlag}) bezeichnen. \textbf{Festlegung:}
Bemessungsgrundlage des gesetzlichen Zuschlags ist der Tarifbeitrag
$b_{x_e}$ ohne Risikozuschlag.

\subsection{Risikozuschläge}\label{sec:risikozuschlag}

Risikozuschläge für erhöhte Risiken (Vorerkrankungen, Beruf) werden als
prozentualer Faktor $RZ$ auf den Bruttobeitrag erhoben. Die Bemessung
im Einzelfall ist nicht Bestandteil dieser TB. Bei Beitragsanpassungen
bleiben die Risikozuschlagsfaktoren unverändert
(Abschnitt~\ref{sec:bap}).
